\documentclass[10pt,a4paper]{article}
\usepackage[T1]{fontenc}
\usepackage{lmodern}
\usepackage[margin=2cm]{geometry}
\usepackage{booktabs}
\usepackage{pgfplots}
\usepackage{pgfplotstable}
\usepgfplotslibrary{groupplots}
\pgfplotsset{compat=1.18}
\usepackage{float}
\usepackage[hidelinks]{hyperref}
\newcommand{\NCells}{43}
\newcommand{\NNorm}{42}
\newcommand{\Seeds}{3}
\newcommand{\NTraces}{15}
\newcommand{\NormTIDE}{0.379}
\newcommand{\NormARC}{0.306}
\newcommand{\NormSthree}{0.334}
\newcommand{\NormLFU}{$-0.257$}
\newcommand{\WinsTIDE}{24}
\newcommand{\GeARC}{34}
\newcommand{\Worst}{$-0.54$}
\newcommand{\WorstCell}{phases, $c=20181$}
\newcommand{\GapShare}{37}
\newcommand{\AdvGain}{+48.6}
\newcommand{\GainLifecycle}{+0.9}
\newcommand{\GainCycle}{+34.5}
\newcommand{\GainPhases}{+10.6}
\newcommand{\PFifty}{15}
\newcommand{\PNineNineMax}{30}
\newcommand{\RpsMin}{45}
\newcommand{\RpsMax}{123}
\newcommand{\TrainPct}{21}

\newcommand{\CurveLifecycle}{36}
\newcommand{\CurveCycle}{1884}
\newcommand{\CurvePhases}{1186}


\definecolor{cLRU}{HTML}{9E9E9E}
\definecolor{cLFU}{HTML}{8D6E63}
\definecolor{cARC}{HTML}{1E88E5}
\definecolor{cS3}{HTML}{43A047}
\definecolor{cSIEVE}{HTML}{FB8C00}
\definecolor{cTIDE}{HTML}{D81B60}
\pgfplotsset{
  every axis/.append style={font=\footnotesize, grid=major, grid style={gray!20}},
  pLRU/.style={cLRU, thick, mark=*, mark size=1.2pt},
  pLFU/.style={cLFU, thick, mark=square*, mark size=1.2pt},
  pARC/.style={cARC, thick, mark=triangle*, mark size=1.6pt},
  pS3/.style={cS3, thick, mark=diamond*, mark size=1.6pt},
  pSIEVE/.style={cSIEVE, thick, mark=pentagon*, mark size=1.4pt},
  pTIDE/.style={cTIDE, very thick, mark=star, mark size=2pt},
  pBEL/.style={black, dashed, thin},
}

\title{TIDE: Trend-Informed Delayed-label Eviction\\\large Benchmark results}
\author{Cache replacement project}
\date{\today}

\begin{document}
\maketitle

\section*{Summary}
TIDE is an ARC cache in which a small neural network, trained online, chooses which item to evict.
ARC's own rule picks the list to evict from. A $12\to64\to1$ MLP, written in pure NumPy, then scores
that list's 4 least-recently-used keys, its most-recently-used key, and the incoming key. Choosing the
incoming key means it is not cached at all (bypass). The network predicts
$\log(1+d/c)$, where $d$ is the time until the key is requested again and $c$ is the cache size.
Every scored key gets its true $d$ as a training label when it is next requested, or a censored label
after $W=8c$ requests, so the model keeps learning as the cache runs. A shadow copy of plain ARC runs
alongside; whenever TIDE falls behind it, TIDE hands decisions back to ARC.

\begin{itemize}
  \item \textbf{Best overall hit rate.} Mean normalized hit rate is \NormTIDE{} for TIDE, against \NormSthree{} for S3-FIFO,
        \NormARC{} for ARC and \NormLFU{} for LFU. Normalized hit rate is $(r-r_{\mathrm{LRU}})/(r_{\mathrm{Belady}}-r_{\mathrm{LRU}})$
        averaged over \NNorm{} cells: 0 means no better than LRU, 1 means as good as the offline optimum.
        TIDE has the highest hit rate in \WinsTIDE{} of \NCells{} cells.
  \item \textbf{Close to ARC's floor.} TIDE $\ge$ ARC in \GeARC{} of \NCells{} cells. Its worst cell is \Worst{} points (\WorstCell).
  \item \textbf{Trend workloads, best gain over ARC:} cycle \GainCycle, phases \GainPhases, lifecycle \GainLifecycle{} points.
        The adversarial loop gains \AdvGain{} points. TIDE realizes \GapShare\% of the gap between ARC and the TIDE oracle on the trend traces.
  \item \textbf{Fast.} A decision takes a median of \PFifty{}\,\textmu s at p50 and at most \PNineNineMax{}\,\textmu s at p99.
        Throughput is \RpsMin k--\RpsMax k requests/s, and hits never call the model.
\end{itemize}

\section*{Setup}
\textbf{Policies.} LRU; LFU (standard in-cache LFU, which forgets counts on eviction); ARC; S3-FIFO; SIEVE; TIDE
(mean over \Seeds{} seeds); the \emph{TIDE oracle}, which has the same design but scores keys with the true next-request time,
giving the design's ceiling; and Belady's MIN, the offline optimum. With bypass allowed, Belady is not a strict upper bound.

\textbf{Traces.} \NTraces{} traces in total. There are 12 existing CSV traces of 100k requests: Zipf, web traffic with a drifting hot set, periodic, bursty,
adversarial, Gaussian, uniform, and Zipf with $\alpha\in\{0.5,\dots,1.5\}$. There are also 3 generated 1M-request trend traces:
\emph{lifecycle}, where items are born and their popularity rises then decays; \emph{cycle}, where 40 groups of 100 keys recur in a fixed
order over a Zipf background; and \emph{phases}, which runs Zipf A, a scan, a loop and Zipf B, twice.
Cache sizes are 0.1\%, 1\% and 10\% of each trace's distinct keys (minimum 16).

\section*{Results}

\begin{figure}[H]
\centering
\begin{tikzpicture}
\pgfplotstableread{gen/norm.dat}\normtab
\begin{axis}[xbar, width=0.55\linewidth, height=5cm, bar width=9pt,
  ytick=data, yticklabels from table={\normtab}{policy}, y dir=reverse,
  xlabel={mean normalized hit rate (0 = LRU, 1 = Belady)}, xmin=-0.42, xmax=0.47,
  nodes near coords, nodes near coords style={font=\tiny, /pgf/number format/fixed, /pgf/number format/precision=3},
  enlarge y limits=0.12]
\addplot[fill=cTIDE!70, draw=cTIDE] table[x=norm, y expr=\coordindex] {\normtab};
\end{axis}
\end{tikzpicture}
\hfill
\begin{minipage}[b]{0.4\linewidth}
\centering\small
\begin{tabular}{lrrr}
\toprule
Policy & Norm. & Wins & Mean hit\% \\
\midrule
TIDE & 0.379 & 24 & 46.09 \\
S3-FIFO & 0.334 & 11 & 44.00 \\
ARC & 0.306 & 3 & 43.36 \\
SIEVE & 0.277 & 2 & 42.92 \\
LRU & 0.000 & 1 & 38.85 \\
LFU & $-$0.257 & 2 & 35.24 \\
\bottomrule
\end{tabular}

\vspace{1.2em}
\end{minipage}
\caption{Overall ranking across all cells. ``Wins'' counts the cells where the policy has the highest hit rate among the six contenders.}
\end{figure}

\begin{figure}[H]
\centering
\begin{tikzpicture}
\pgfplotstableread{gen/delta.dat}\deltatab
\begin{axis}[ybar=0pt, bar width=3pt, clip=false, width=\linewidth, height=6cm,
  xtick=data, xticklabels from table={\deltatab}{trace}, xticklabel style={rotate=45, anchor=east, font=\scriptsize},
  ylabel={TIDE $-$ ARC (hit-rate points)}, ymin=-1.5, ymax=12, enlarge x limits=0.03,
  legend style={at={(0.99,0.97)}, anchor=north east, font=\scriptsize}, legend columns=3,
  extra y ticks={0}, extra y tick style={grid style={black, thin}}]
\addplot[fill=cARC!35, draw=cARC!60] table[x expr=\coordindex, y=s1] {\deltatab};
\addplot[fill=cARC!70, draw=cARC] table[x expr=\coordindex, y=s2] {\deltatab};
\addplot[fill=cTIDE!80, draw=cTIDE] table[x expr=\coordindex, y=s3] {\deltatab};
\legend{0.1\% cache, 1\% cache, 10\% cache}
\node[above,font=\tiny,inner sep=1pt] at ([xshift=3pt]axis cs:1,12) {+34};
\node[above,font=\tiny,inner sep=1pt] at ([xshift=0pt]axis cs:5,12) {+49};
\end{axis}
\end{tikzpicture}
\caption{Gain of TIDE over ARC per trace and cache size. Bars above 12 points are clipped and labelled with their value.}
\end{figure}

\begin{figure}[H]
\centering
\begin{tikzpicture}
\begin{groupplot}[group style={group size=3 by 1, horizontal sep=1.2cm}, width=0.31\linewidth, height=5cm,
  xmode=log, xlabel={cache size (\% of keys)}, xtick={0.1,1,10}, xticklabels={0.1,1,10}]
\nextgroupplot[title=lifecycle, ylabel={hit rate (\%)}]
\addplot[pLRU] table[x=frac, y=LRU] {gen/size_lifecycle.dat};
\addplot[pLFU] table[x=frac, y=LFU] {gen/size_lifecycle.dat};
\addplot[pARC] table[x=frac, y=ARC] {gen/size_lifecycle.dat};
\addplot[pS3] table[x=frac, y=S3FIFO] {gen/size_lifecycle.dat};
\addplot[pSIEVE] table[x=frac, y=SIEVE] {gen/size_lifecycle.dat};
\addplot[pTIDE] table[x=frac, y=TIDE] {gen/size_lifecycle.dat};
\addplot[pBEL] table[x=frac, y=Belady] {gen/size_lifecycle.dat};
\nextgroupplot[title=cycle]
\addplot[pLRU] table[x=frac, y=LRU] {gen/size_cycle.dat};
\addplot[pLFU] table[x=frac, y=LFU] {gen/size_cycle.dat};
\addplot[pARC] table[x=frac, y=ARC] {gen/size_cycle.dat};
\addplot[pS3] table[x=frac, y=S3FIFO] {gen/size_cycle.dat};
\addplot[pSIEVE] table[x=frac, y=SIEVE] {gen/size_cycle.dat};
\addplot[pTIDE] table[x=frac, y=TIDE] {gen/size_cycle.dat};
\addplot[pBEL] table[x=frac, y=Belady] {gen/size_cycle.dat};
\nextgroupplot[title=phases, legend style={at={(1.05,0.5)}, anchor=west, font=\scriptsize}]
\addplot[pLRU] table[x=frac, y=LRU] {gen/size_phases.dat};
\addplot[pLFU] table[x=frac, y=LFU] {gen/size_phases.dat};
\addplot[pARC] table[x=frac, y=ARC] {gen/size_phases.dat};
\addplot[pS3] table[x=frac, y=S3FIFO] {gen/size_phases.dat};
\addplot[pSIEVE] table[x=frac, y=SIEVE] {gen/size_phases.dat};
\addplot[pTIDE] table[x=frac, y=TIDE] {gen/size_phases.dat};
\addplot[pBEL] table[x=frac, y=Belady] {gen/size_phases.dat};
\legend{LRU, LFU, ARC, S3-FIFO, SIEVE, TIDE, Belady}
\end{groupplot}
\end{tikzpicture}
\caption{Hit rate vs.\ cache size on the three 1M-request trend traces.}
\end{figure}

\begin{figure}[H]
\centering
\begin{tikzpicture}
\begin{groupplot}[group style={group size=3 by 1, horizontal sep=1.4cm}, width=0.36\linewidth, height=5cm,
  xlabel={position in trace (\%)}, xmin=0, xmax=100, no markers]
\nextgroupplot[title={lifecycle, $c=\CurveLifecycle$}, ylabel={hit rate per 1\% window (\%)}]
\addplot[cARC, thick] table[x=pct, y=ARC] {gen/curve_lifecycle.dat};
\addplot[cTIDE, thick] table[x=pct, y=TIDE] {gen/curve_lifecycle.dat};
\nextgroupplot[title={cycle, $c=\CurveCycle$}]
\addplot[cARC, thick] table[x=pct, y=ARC] {gen/curve_cycle.dat};
\addplot[cTIDE, thick] table[x=pct, y=TIDE] {gen/curve_cycle.dat};
\nextgroupplot[title={phases, $c=\CurvePhases$}, legend style={at={(0.98,0.98)}, anchor=north east, font=\scriptsize}]
\addplot[cARC, thick] table[x=pct, y=ARC] {gen/curve_phases.dat};
\addplot[cTIDE, thick] table[x=pct, y=TIDE] {gen/curve_phases.dat};
\legend{ARC, TIDE}
\end{groupplot}
\end{tikzpicture}
\caption{Learning live: hit rate over time for ARC and TIDE, at the cache size where TIDE gains most on each trend trace.
TIDE starts out as plain ARC and switches to its model once it has $\ge$\,2048 labels and $t\ge W$.}
\end{figure}

\begin{table}[H]
\centering\small
\begin{tabular}{llrrrr}
\toprule
Trace & Implementation & Req/s & p50 ($\mu$s) & p99 ($\mu$s) & p99.9 ($\mu$s) \\
\midrule
MSR hm\_0 & ARC & 2,487k & 0.16 & 1.0 & 3.1 \\
 & TIDE, NumPy & 128k & 2.56 & 133.5 & 183.3 \\
 & TIDE, Numba sync & 150k & 2.49 & 87.3 & 134.3 \\
 & TIDE, Numba async & 155k & 2.90 & 25.8 & 64.1 \\
\midrule
Twitter c50 & ARC & 1,099k & 0.79 & 1.8 & 4.3 \\
 & TIDE, NumPy & 48k & 14.15 & 219.3 & 410.8 \\
 & TIDE, Numba sync & 83k & 9.62 & 112.2 & 125.2 \\
 & TIDE, Numba async & 108k & 9.43 & 20.8 & 58.2 \\
\midrule
Meta CDN reag & ARC & 1,649k & 0.28 & 1.0 & 5.0 \\
 & TIDE, NumPy & 81k & 7.16 & 186.8 & 221.9 \\
 & TIDE, Numba sync & 122k & 4.99 & 95.9 & 110.6 \\
 & TIDE, Numba async & 145k & 5.12 & 17.5 & 62.7 \\
\bottomrule
\end{tabular}

\caption{Speed, measured serially at the 1\% cache size in pure Python/NumPy. Latency is per eviction decision and covers
features, label snapshot and inference; training steps are excluded. Training \% is the share of total runtime spent in
online training.}
\end{table}

\section*{Caveats}
\begin{itemize}
  \item LFU here is the standard in-cache LFU. An LFU with global, never-forgotten counts (unbounded metadata) beats TIDE on the i.i.d.\ Zipf traces;
        it is in \texttt{bench\_results.csv} as \texttt{LFU*} and is left out of the contenders.
  \item The trend generators and their parameters were written for this project and fixed before tuning; tuning used the same settings on every trace.
        TIDE's settings ($\delta=0.1$ for tail keys, $0.2$ for MRU/bypass, one training step per 32 labels, 64 hidden units) come from a single sweep over all traces.
  \item Pure Python runs 20--70$\times$ below plain ARC's raw throughput; a Numba or C hot path would close that.
  \item TIDE missed some internal targets: worst cell $-0.54$ (target $\ge-0.5$), lifecycle gain $<1$ point, and training at about \TrainPct\% of runtime (target $\le 15\%$).
\end{itemize}

\clearpage
\begin{table}[H]
\centering\footnotesize
\begin{tabular}{lrrrrrrrrrrrr}
\toprule
Trace & Size\% & $c$ & LRU & ARC & S3-FIFO & LIRS & Cacheus & LRB & TIDE & Oracle & Belady & $\Delta$ARC \\
\midrule
adversarial & 10 & 16 & 0.00 & 0.00 & 0.00 & 48.38 & 0.00 & 0.00 & \textbf{48.55\,$\pm$\,1.94} & 51.60 & 49.99 & +48.55 \\
\midrule
bursty & 0.1 & 16 & 84.32 & 84.32 & 84.32 & 84.32 & \textbf{84.33} & 84.32 & 84.25\,$\pm$\,0.17 & 84.76 & 85.04 & $-$0.08 \\
 & 1 & 80 & \textbf{84.43} & 84.42 & 84.42 & 84.42 & 84.43 & \textbf{84.43} & 84.26\,$\pm$\,0.18 & 85.22 & 86.08 & $-$0.16 \\
 & 10 & 801 & \textbf{85.52} & 85.51 & 85.50 & 85.51 & 85.47 & \textbf{85.52} & 85.19\,$\pm$\,0.24 & 87.04 & 89.36 & $-$0.32 \\
\midrule
cycle & 0.1 & 19 & 6.30 & 16.03 & \textbf{16.32} & 16.03 & 14.41 & 12.37 & 16.23\,$\pm$\,0.06 & 16.93 & 17.57 & +0.19 \\
 & 1 & 188 & 15.03 & 22.05 & 21.92 & \textbf{22.22} & 21.26 & 19.11 & 21.93\,$\pm$\,0.25 & 22.25 & 23.72 & $-$0.12 \\
 & 10 & 1884 & 20.87 & 24.29 & 24.15 & \textbf{50.36} & 41.57 & 37.73 & 42.71\,$\pm$\,0.88 & 43.71 & 53.21 & +18.43 \\
\midrule
gaussian & 0.1 & 16 & 0.26 & 0.44 & \textbf{0.48} & 0.45 & 0.27 & 0.26 & 0.41\,$\pm$\,0.08 & 5.12 & 6.22 & $-$0.03 \\
 & 1 & 88 & 1.50 & 1.97 & \textbf{2.05} & 2.01 & 1.72 & 1.50 & 1.94\,$\pm$\,0.05 & 10.86 & 15.88 & $-$0.03 \\
 & 10 & 880 & 14.78 & 16.79 & 17.15 & 16.99 & 16.48 & 14.78 & \textbf{17.97\,$\pm$\,0.10} & 32.55 & 47.55 & +1.18 \\
\midrule
lifecycle & 0.1 & 36 & 33.14 & 42.74 & 43.90 & 43.32 & 43.07 & 33.14 & \textbf{44.17\,$\pm$\,0.93} & 48.68 & 53.81 & +1.43 \\
 & 1 & 359 & 66.61 & 69.19 & 68.81 & 68.85 & 68.27 & 66.61 & \textbf{69.22\,$\pm$\,0.44} & 73.28 & 81.27 & +0.04 \\
 & 10 & 3592 & 91.41 & \textbf{91.44} & 88.72 & 86.54 & 90.59 & 91.41 & 91.31\,$\pm$\,0.20 & 91.46 & 92.75 & $-$0.13 \\
\midrule
periodic & 0.1 & 16 & 10.27 & 11.64 & 12.39 & 12.45 & 11.28 & 10.27 & \textbf{12.62\,$\pm$\,0.10} & 32.80 & 36.66 & +0.97 \\
 & 1 & 89 & 51.00 & 66.80 & 63.79 & \textbf{67.84} & 67.03 & 51.00 & 67.31\,$\pm$\,0.39 & 72.11 & 76.36 & +0.51 \\
 & 10 & 891 & \textbf{80.11} & 80.03 & 79.60 & 78.54 & 79.88 & \textbf{80.11} & 79.73\,$\pm$\,0.17 & 81.94 & 85.56 & $-$0.30 \\
\midrule
phases & 0.1 & 119 & 19.17 & 26.17 & 26.41 & 26.30 & \textbf{27.88} & 20.23 & 27.16\,$\pm$\,0.65 & 29.20 & 32.26 & +1.00 \\
 & 1 & 1186 & 32.57 & 36.77 & 36.50 & 36.79 & \textbf{46.58} & 33.21 & 40.87\,$\pm$\,3.60 & 53.12 & 60.44 & +4.09 \\
 & 10 & 11861 & 69.42 & 69.26 & \textbf{69.61} & 68.45 & 69.13 & 69.42 & 69.01\,$\pm$\,0.07 & 70.21 & 74.50 & $-$0.25 \\
\midrule
uniform & 0.1 & 16 & 0.15 & 0.17 & 0.17 & \textbf{0.17} & 0.14 & 0.15 & 0.16\,$\pm$\,0.02 & 3.75 & 4.82 & $-$0.01 \\
 & 1 & 100 & 1.01 & 1.00 & 1.01 & \textbf{1.03} & 0.98 & 1.01 & 1.02\,$\pm$\,0.04 & 8.56 & 13.17 & +0.02 \\
 & 10 & 1000 & 10.04 & \textbf{10.14} & 10.05 & 10.07 & 9.96 & 10.04 & 10.04\,$\pm$\,0.10 & 25.42 & 40.05 & $-$0.09 \\
\midrule
web & 0.1 & 16 & 3.39 & 4.52 & \textbf{5.18} & 4.88 & 3.93 & 3.39 & 4.86\,$\pm$\,0.09 & 18.73 & 21.59 & +0.34 \\
 & 1 & 50 & 10.49 & 14.18 & 14.87 & 14.89 & 12.80 & 10.49 & \textbf{14.96\,$\pm$\,0.04} & 31.43 & 38.99 & +0.77 \\
 & 10 & 499 & 67.74 & \textbf{71.35} & 70.79 & 70.73 & 70.44 & 67.74 & 71.04\,$\pm$\,0.07 & 73.04 & 78.85 & $-$0.32 \\
\midrule
zipf-100k & 0.1 & 19 & 35.46 & 47.67 & \textbf{49.27} & 48.18 & 47.59 & 35.46 & 49.25\,$\pm$\,0.05 & 51.22 & 53.53 & +1.58 \\
 & 1 & 194 & 59.43 & 67.28 & 67.35 & 67.45 & 67.48 & 59.43 & \textbf{67.57\,$\pm$\,0.03} & 68.19 & 71.26 & +0.29 \\
 & 10 & 1935 & 74.08 & 77.00 & 76.65 & 77.15 & 76.85 & 74.08 & \textbf{77.24\,$\pm$\,0.13} & 78.22 & 80.16 & +0.24 \\
\midrule
zipf-a0.5 & 0.1 & 16 & 0.38 & 2.20 & \textbf{2.66} & 2.22 & 0.68 & 0.38 & 2.36\,$\pm$\,0.09 & 6.26 & 7.20 & +0.16 \\
 & 1 & 100 & 2.39 & 6.58 & \textbf{7.35} & 6.70 & 3.15 & 2.39 & 7.20\,$\pm$\,0.09 & 13.04 & 18.26 & +0.61 \\
 & 10 & 999 & 17.98 & 23.67 & 25.05 & 23.80 & 22.76 & 17.98 & \textbf{25.77\,$\pm$\,0.10} & 32.38 & 48.12 & +2.10 \\
\midrule
zipf-a0.8 & 0.1 & 16 & 4.29 & 13.63 & \textbf{14.87} & 13.79 & 12.02 & 4.29 & 14.53\,$\pm$\,0.22 & 17.92 & 19.90 & +0.91 \\
 & 1 & 97 & 15.50 & 26.21 & 27.14 & 26.48 & 25.94 & 15.50 & \textbf{27.35\,$\pm$\,0.20} & 30.18 & 36.63 & +1.14 \\
 & 10 & 970 & 43.06 & 50.34 & 51.46 & 50.41 & 51.19 & 43.06 & \textbf{51.82\,$\pm$\,0.12} & 54.17 & 65.75 & +1.48 \\
\midrule
zipf-a1.0 & 0.1 & 16 & 17.98 & 30.75 & 32.75 & 31.20 & 30.15 & 17.98 & \textbf{32.75\,$\pm$\,0.15} & 35.77 & 38.35 & +2.00 \\
 & 1 & 85 & 37.12 & 48.12 & 48.97 & 48.45 & 48.65 & 37.12 & \textbf{49.27\,$\pm$\,0.11} & 51.04 & 56.49 & +1.15 \\
 & 10 & 853 & 65.33 & 70.72 & 71.27 & 70.77 & \textbf{71.47} & 65.33 & 71.43\,$\pm$\,0.09 & 72.20 & 79.40 & +0.71 \\
\midrule
zipf-a1.2 & 0.1 & 16 & 33.30 & 45.70 & 47.51 & 46.18 & 45.69 & 33.30 & \textbf{47.57\,$\pm$\,0.09} & 49.72 & 51.83 & +1.88 \\
 & 1 & 89 & 52.82 & 61.74 & 62.30 & 61.96 & 62.19 & 52.82 & \textbf{62.42\,$\pm$\,0.11} & 63.31 & 67.02 & +0.68 \\
 & 10 & 886 & 72.01 & 75.57 & \textbf{75.84} & 75.64 & 75.77 & 72.01 & 75.83\,$\pm$\,0.10 & 76.88 & 82.21 & +0.26 \\
\midrule
zipf-a1.5 & 0.1 & 16 & 72.58 & 77.96 & 79.82 & 78.63 & 79.05 & 72.58 & \textbf{80.02\,$\pm$\,0.14} & 81.21 & 82.89 & +2.06 \\
 & 1 & 28 & 79.19 & 83.40 & 84.46 & 83.86 & 84.29 & 79.19 & \textbf{84.61\,$\pm$\,0.08} & 85.48 & 87.28 & +1.21 \\
 & 10 & 282 & 93.37 & 94.68 & 94.69 & 94.71 & 94.74 & 93.37 & \textbf{94.75\,$\pm$\,0.02} & 94.99 & 95.95 & +0.07 \\
\bottomrule
\end{tabular}

\caption{Hit rate (\%) for every cell. Bold marks the best of the six contenders. Oracle is the TIDE oracle (true next-request scores);
$\Delta$ARC is TIDE $-$ ARC. TIDE is the mean over \Seeds{} seeds.}
\end{table}

\end{document}
